\needspace{5cm} \section{Improvizační trojzubec} \label{improvizační trojzubec} Pojem označující seřazenou sadu tří základních otázek, na které by si měla umět každá \odkaz{postava}{postava} i \odkaz{hráč}{hráč} ji stvárňující v~improvizaci odpovědět. Pokud to nedokáže, měla by ve scénce vytvořit věci, které k~tomu chybí. 
 
\subsubsection{ Kdo jste? } Byly všechny postavy dobře popsány? Ví všichni o~všech dost na to, aby z~toho mohli vycházet? 
 
\subsubsection{ Kde jste? } Bylo definováno \odkaz{prostředí}{prostředí}? Dá se s~ním pracovat? 
 
\subsubsection{ Co tam děláte? } Proč jsme v~této situaci? Proč moje postava dělá to, co dělá? Zachovala by se takhle pokaždé? 
 
\subsection{ Viz také } Mezi podobné pomocné techniky patří \odkaz{CROW}{crow} nebo třeba \odkaz{5 otázek}{5 otázek}. 
 
 
 
\needspace{5cm} \section{5 otázek} \label{5 otázek} 5 otázek je technika nápomocných k~vytváření scény. V~čechách ji učili např. David Weber či Jim Libby i s~použitím mnemotechnických gest. 
Jádrem techniky je, že si hráč zodpoví několik otázek. 
 
Je důležité zmínit, že na otázky nemusíme odpovídat nutně v~tomto pořadí, většinou zjištění v~jedné oblasti implikuje či nabízí odpovědi na další otázky. 
 
\subsubsection{ Kdo jsem já? } Co vím o~své \odkaz{postavě}{postava}. Jaké má cíle/chování/vlastnosti. 
\subsubsection{ Kdo jsi ty? } Co vím o~postavě dalšího hráče/hráčů. 
\subsubsection{ Kdo jsme my, jaký je mezi námi vztah? } Jaký je \odkaz{vztah}{vztah} mezi našimi postavami?  
\subsubsection{ Kde jsme? } V~jakém prostředí se nacházíme, jaké emoce v~nás samotné prostředí vyvolává, kde je rozmístěný jaký nábytek/součásti scény. 
\subsubsection{ Co tady děláme? } Jaký je cíl naší scény, jaká je hra scény, co se o~našich postavách dalšího divák dozvídá. 
 
 
 
K~těmto 5 otázkám má smysl u~některých \odkaz{forem}{forma} improvizace si odpovědět další pomocné otázky, např. V~jakém \odkaz{žánru}{žánr} hrajeme? Kdy se scéna vzhledem k~již sehraným odehrává časově? V~jaké fázi \odkaz{příběhu}{příběh} se nacházíme. 
 
\subsection{ Viz také } \begin{itemize}
\item  \odkaz{CROW}{crow}
\item  \odkaz{Improvizační trojzubec}{improvizační trojzubec}
\end{itemize}
 
 
 
\needspace{5cm} \section{CROW} \label{crow} CROW je technika otázek nápomocných  vytváření reality scény v~improvizaci. Název je zkratkou. 
 
\subsubsection{ C - Character - Postava } Co o~\odkaz{postavě}{postava} víme? Jaké má vlastnosti, vzhled, chování?  
 
\subsubsection{ R- Relation - Vztah } Jaké jsou \odkaz{vztahy}{vztah} mezi postavami, jak na ostatní protagonista nahlíží. 
 
\subsubsection{ O~- Objective - Cíl } Po čem postavy / příběh touží a jaké překážky v~jejich cestě leží. 
 
\subsubsection{ W- Where/When - Kde a Kdy } Kdy a kde se scéna odehrává? V~jakém \odkaz{prostředí}{prostředí}? Ať již prostorově, časově, nebo ve vztahu k~předešlým scénám. 
 
 
\subsection{ Viz také } \begin{itemize}
\item  \odkaz{Improvizační trojzubec}{improvizační trojzubec}
\item  \odkaz{5 otázek}{5 otázek}
\end{itemize}
 
 
 
